% ============================================================
% Feature Matrix Table — From Raw Data to Features
% CSC491 Textbook, Chapter: Machine Learning
%
% Each row = one dollar bar (a moment in market time)
% Columns = engineered features (x) + label (y)
%
% Preamble requirements:
%     \usepackage{booktabs, array, xcolor}
% ============================================================

\begin{table}[ht]
\centering
\setlength{\tabcolsep}{10pt}
\renewcommand{\arraystretch}{1.6}
\caption{%
  A financial feature matrix: each row represents one dollar volume bar,
  and each column represents a measurable signal at that moment in time.
  The columns are \textbf{features} $x$
}
\label{tab:feature-matrix}
\begin{tabular}{
  >{\centering\arraybackslash}p{2.0cm}   % timestamp
  >{\centering\arraybackslash}p{1.8cm}   % close
  >{\centering\arraybackslash}p{1.8cm}   % EMA
  >{\centering\arraybackslash}p{1.8cm}   % RSI
  >{\centering\arraybackslash}p{1.8cm}   % ATR
}
\toprule
\multicolumn{5}{c}{\textbf{Features} $x$} \\
\cmidrule(r){1-5}
\textbf{Timestamp}
  & \textbf{Close}
  & \textbf{EMA\_14}
  & \textbf{RSI\_14}
  & \textbf{ATR\_14} \\
\midrule
09:32:01 & 148.20 & 147.80 & 54.2 & 1.32 \\
09:34:47 & 149.50 & 148.10 & 61.7 & 1.41 \\
09:38:03 & 148.90 & 148.30 & 58.4 & 1.38 \\
\multicolumn{5}{c}{$\vdots$} \\
\bottomrule
\end{tabular}

\smallskip
\footnotesize
\textit{Note:} In practice the feature matrix will contain hundreds of rows (one per
dollar bar) and potentially dozens of columns.

\end{table}